% GENERATED FILE. Edit canonical lessons, then run npm run generate:books.
% canonical-source-hash: fnv1a64:6f20bccc3c3a9086
% canonical-lessons: PT-C310-batizado, PT-C310-batizar, PT-C310-padrinho, PT-C310-madrinha, PT-C310-amamentar

\chapter{O batizado}
\label{ch:pt-o-batizado}
\hlchaptermodality{310}

\begin{chapteropening}
\textbf{By the end of this chapter:} \emph{I can talk about a christening, godparents, and the first months of a baby’s life.}

Complete the last lesson of chapter 310: I can talk about a christening, godparents, and the first months of a baby’s life.
\end{chapteropening}

\section[o batizado]{o batizado — the christening, the baptism}

\label{lesson:PT-C310-batizado}



\begin{quote}
Before the new one: say the Portuguese for pregnancy, then the Portuguese for the newborn baby.
\end{quote}

\subsection*{You'll want to know: o batizado}
\textbf{o batizado} — \textquotedblleft{}the christening, the baptism\textquotedblright{}.

\begin{quote}
O batizado foi na igreja da aldeia.
\end{quote}

\begin{quote}
\emph{The christening was in the village church.}
\end{quote}

Before the 1990 spelling reform, \emph{baptizado}.

\begin{cousinweb}
\emph{batizado} is the past participle of \emph{batizar} used as a noun; the verb goes back to Latin \emph{baptizāre}, from Greek \emph{baptízein} \textquotedblleft{}to dip, to plunge\textquotedblright{} (Lewis \& Short).
\end{cousinweb}

\begin{grammarlens}[title={What you've built}]
One more word for telling, explaining and describing at B1.
\end{grammarlens}

\subsection*{Your turn}
\begin{itemize}
\raggedright
  \item \emph{Say it:} \emph{o batizado}
  \item \emph{Say it:} \emph{o batizado}, once more
  \item \emph{Say it:} a sentence of your own with \emph{o batizado}
  \item \emph{Recall:} say the Portuguese for pregnancy, then the Portuguese for the newborn baby, then say \emph{o batizado} again
\end{itemize}

\begin{tcolorbox}[breakable,colback=teal!4,colframe=teal!35!black,title={Before you move on}]
What does o batizado mean? (\textquotedblleft{}The christening, the baptism\textquotedblright{}.) Say it once more.
\end{tcolorbox}
\section[batizar]{batizar — to baptise, to christen; to name}

\label{lesson:PT-C310-batizar}



\begin{quote}
Before the new one: say the Portuguese for the newborn baby, then the Portuguese for the christening.
\end{quote}

\subsection*{You'll want to know: batizar}
\textbf{batizar} — \textquotedblleft{}to baptise, to christen; to name\textquotedblright{}.

\begin{quote}
Batizaram a menina em junho.
\end{quote}

\begin{quote}
\emph{They christened the little girl in June.}
\end{quote}

Also to give a name to anything: \emph{batizaram o barco com o nome da avó}.

\begin{cousinweb}
\emph{batizar} comes down from Old Portuguese \emph{bautizar} (EtymDB), from Latin \emph{baptizāre}, from Greek \emph{baptízein} \textquotedblleft{}to dip\textquotedblright{}.
\end{cousinweb}

\begin{grammarlens}[title={What you've built}]
One more word for telling, explaining and describing at B1.
\end{grammarlens}

\subsection*{Your turn}
\begin{itemize}
\raggedright
  \item \emph{Say it:} \emph{batizar}
  \item \emph{Say it:} \emph{batizar}, once more
  \item \emph{Say it:} a sentence of your own with \emph{batizar}
  \item \emph{Recall:} say the Portuguese for the newborn baby, then the Portuguese for the christening, then say \emph{batizar} again
\end{itemize}

\begin{tcolorbox}[breakable,colback=teal!4,colframe=teal!35!black,title={Before you move on}]
What does batizar mean? (\textquotedblleft{}To baptise, to christen; to name\textquotedblright{}.) Say it once more.
\end{tcolorbox}
\section[o padrinho]{o padrinho — the godfather; the best man}

\label{lesson:PT-C310-padrinho}



\begin{quote}
Before the new one: say the Portuguese for the christening, then the Portuguese for to baptise.
\end{quote}

\subsection*{You'll want to know: o padrinho}
\textbf{o padrinho} — \textquotedblleft{}the godfather; the best man\textquotedblright{}.

\begin{quote}
O meu tio é o meu padrinho.
\end{quote}

\begin{quote}
\emph{My uncle is my godfather.}
\end{quote}

At a wedding, \emph{os padrinhos} are the witnesses.

\begin{cousinweb}
\emph{padrinho} goes back to Vulgar Latin \emph{patrīnus} \textquotedblleft{}godfather\textquotedblright{} (EtymDB), from \emph{pater} \textquotedblleft{}father\textquotedblright{}; Italian and Spanish \emph{padrino} are its cousins.
\end{cousinweb}

\begin{grammarlens}[title={What you've built}]
One more word for telling, explaining and describing at B1.
\end{grammarlens}

\subsection*{Your turn}
\begin{itemize}
\raggedright
  \item \emph{Say it:} \emph{o padrinho}
  \item \emph{Say it:} \emph{o padrinho}, once more
  \item \emph{Say it:} a sentence of your own with \emph{o padrinho}
  \item \emph{Recall:} say the Portuguese for the christening, then the Portuguese for to baptise, then say \emph{o padrinho} again
\end{itemize}

\begin{tcolorbox}[breakable,colback=teal!4,colframe=teal!35!black,title={Before you move on}]
What does o padrinho mean? (\textquotedblleft{}The godfather; the best man\textquotedblright{}.) Say it once more.
\end{tcolorbox}
\section[a madrinha]{a madrinha — the godmother}

\label{lesson:PT-C310-madrinha}



\begin{quote}
Before the new one: say the Portuguese for to baptise, then the Portuguese for the godfather.
\end{quote}

\subsection*{You'll want to know: a madrinha}
\textbf{a madrinha} — \textquotedblleft{}the godmother\textquotedblright{}.

\begin{quote}
A madrinha deu-lhe um presente.
\end{quote}

\begin{quote}
\emph{Her godmother gave her a present.}
\end{quote}

\emph{O afilhado}, \emph{a afilhada} — the godchild.

\begin{cousinweb}
\emph{madrinha} comes down from Old Portuguese \emph{madrỹa} (EtymDB), from Vulgar Latin \emph{matrīna}, built on Latin \emph{māter} \textquotedblleft{}mother\textquotedblright{}.
\end{cousinweb}

\begin{grammarlens}[title={What you've built}]
One more word for telling, explaining and describing at B1.
\end{grammarlens}

\subsection*{Your turn}
\begin{itemize}
\raggedright
  \item \emph{Say it:} \emph{a madrinha}
  \item \emph{Say it:} \emph{a madrinha}, once more
  \item \emph{Say it:} a sentence of your own with \emph{a madrinha}
  \item \emph{Recall:} say the Portuguese for to baptise, then the Portuguese for the godfather, then say \emph{a madrinha} again
\end{itemize}

\begin{tcolorbox}[breakable,colback=teal!4,colframe=teal!35!black,title={Before you move on}]
What does a madrinha mean? (\textquotedblleft{}The godmother\textquotedblright{}.) Say it once more.
\end{tcolorbox}
\section[amamentar]{amamentar — to breastfeed}

\label{lesson:PT-C310-amamentar}



\begin{quote}
Before the new one: say the Portuguese for the godfather, then the Portuguese for the godmother.
\end{quote}

\subsection*{You'll want to know: amamentar}
\textbf{amamentar} — \textquotedblleft{}to breastfeed\textquotedblright{}.

\begin{quote}
Ela amamentou o filho durante seis meses.
\end{quote}

\begin{quote}
\emph{She breastfed her son for six months.}
\end{quote}

Built on \emph{a mama}, ‘the breast’.

\begin{cousinweb}
\emph{amamentar} is built from \emph{a-} and \emph{mama} \textquotedblleft{}breast\textquotedblright{} (EtymDB), from Latin \emph{mamma} \textquotedblleft{}breast\textquotedblright{}.
\end{cousinweb}

\begin{grammarlens}[title={What you've built}]
One more word for telling, explaining and describing at B1.
\end{grammarlens}

\subsection*{Your turn}
\begin{itemize}
\raggedright
  \item \emph{Say it:} \emph{amamentar}
  \item \emph{Say it:} \emph{amamentar}, once more
  \item \emph{Say it:} a sentence of your own with \emph{amamentar}
  \item \emph{Recall:} say the Portuguese for the godfather, then the Portuguese for the godmother, then say \emph{amamentar} again
\end{itemize}

\begin{tcolorbox}[breakable,colback=teal!4,colframe=teal!35!black,title={Before you move on}]
What does amamentar mean? (\textquotedblleft{}To breastfeed\textquotedblright{}.) Say it once more.
\end{tcolorbox}
